\IfFileExists{stacks-project.cls}{%
\documentclass{stacks-project}
}{%
\documentclass{amsart}
}

% For dealing with references we use the comment environment
\usepackage{verbatim}
\newenvironment{reference}{\comment}{\endcomment}
%\newenvironment{reference}{}{}
\newenvironment{slogan}{\comment}{\endcomment}
\newenvironment{history}{\comment}{\endcomment}

% For commutative diagrams we use Xy-pic
\usepackage[all]{xy}

% We use 2cell for 2-commutative diagrams.
\xyoption{2cell}
\UseAllTwocells

% We use multicol for the list of chapters between chapters
\usepackage{multicol}

% This is generally recommended for better output
\usepackage{lmodern}
\usepackage[T1]{fontenc}

% For cross-file-references
\usepackage{xr-hyper}

% Package for hypertext links:
\usepackage{hyperref}

% For any local file, say "hello.tex" you want to link to please
% use \externaldocument[hello-]{hello}
\externaldocument[introduction-]{introduction}
\externaldocument[conventions-]{conventions}
\externaldocument[sets-]{sets}
\externaldocument[categories-]{categories}
\externaldocument[topology-]{topology}
\externaldocument[sheaves-]{sheaves}
\externaldocument[sites-]{sites}
\externaldocument[stacks-]{stacks}
\externaldocument[fields-]{fields}
\externaldocument[algebra-]{algebra}
\externaldocument[brauer-]{brauer}
\externaldocument[homology-]{homology}
\externaldocument[derived-]{derived}
\externaldocument[simplicial-]{simplicial}
\externaldocument[more-algebra-]{more-algebra}
\externaldocument[smoothing-]{smoothing}
\externaldocument[modules-]{modules}
\externaldocument[sites-modules-]{sites-modules}
\externaldocument[injectives-]{injectives}
\externaldocument[cohomology-]{cohomology}
\externaldocument[sites-cohomology-]{sites-cohomology}
\externaldocument[dga-]{dga}
\externaldocument[dpa-]{dpa}
\externaldocument[sdga-]{sdga}
\externaldocument[hypercovering-]{hypercovering}
\externaldocument[schemes-]{schemes}
\externaldocument[constructions-]{constructions}
\externaldocument[properties-]{properties}
\externaldocument[morphisms-]{morphisms}
\externaldocument[coherent-]{coherent}
\externaldocument[divisors-]{divisors}
\externaldocument[limits-]{limits}
\externaldocument[varieties-]{varieties}
\externaldocument[topologies-]{topologies}
\externaldocument[descent-]{descent}
\externaldocument[perfect-]{perfect}
\externaldocument[more-morphisms-]{more-morphisms}
\externaldocument[flat-]{flat}
\externaldocument[groupoids-]{groupoids}
\externaldocument[more-groupoids-]{more-groupoids}
\externaldocument[etale-]{etale}
\externaldocument[chow-]{chow}
\externaldocument[intersection-]{intersection}
\externaldocument[pic-]{pic}
\externaldocument[weil-]{weil}
\externaldocument[adequate-]{adequate}
\externaldocument[dualizing-]{dualizing}
\externaldocument[duality-]{duality}
\externaldocument[discriminant-]{discriminant}
\externaldocument[derham-]{derham}
\externaldocument[local-cohomology-]{local-cohomology}
\externaldocument[algebraization-]{algebraization}
\externaldocument[curves-]{curves}
\externaldocument[resolve-]{resolve}
\externaldocument[models-]{models}
\externaldocument[functors-]{functors}
\externaldocument[equiv-]{equiv}
\externaldocument[pione-]{pione}
\externaldocument[etale-cohomology-]{etale-cohomology}
\externaldocument[proetale-]{proetale}
\externaldocument[relative-cycles-]{relative-cycles}
\externaldocument[more-etale-]{more-etale}
\externaldocument[trace-]{trace}
\externaldocument[crystalline-]{crystalline}
\externaldocument[spaces-]{spaces}
\externaldocument[spaces-properties-]{spaces-properties}
\externaldocument[spaces-morphisms-]{spaces-morphisms}
\externaldocument[decent-spaces-]{decent-spaces}
\externaldocument[spaces-cohomology-]{spaces-cohomology}
\externaldocument[spaces-limits-]{spaces-limits}
\externaldocument[spaces-divisors-]{spaces-divisors}
\externaldocument[spaces-over-fields-]{spaces-over-fields}
\externaldocument[spaces-topologies-]{spaces-topologies}
\externaldocument[spaces-descent-]{spaces-descent}
\externaldocument[spaces-perfect-]{spaces-perfect}
\externaldocument[spaces-more-morphisms-]{spaces-more-morphisms}
\externaldocument[spaces-flat-]{spaces-flat}
\externaldocument[spaces-groupoids-]{spaces-groupoids}
\externaldocument[spaces-more-groupoids-]{spaces-more-groupoids}
\externaldocument[bootstrap-]{bootstrap}
\externaldocument[spaces-pushouts-]{spaces-pushouts}
\externaldocument[spaces-chow-]{spaces-chow}
\externaldocument[groupoids-quotients-]{groupoids-quotients}
\externaldocument[spaces-more-cohomology-]{spaces-more-cohomology}
\externaldocument[spaces-simplicial-]{spaces-simplicial}
\externaldocument[spaces-duality-]{spaces-duality}
\externaldocument[formal-spaces-]{formal-spaces}
\externaldocument[restricted-]{restricted}
\externaldocument[spaces-resolve-]{spaces-resolve}
\externaldocument[formal-defos-]{formal-defos}
\externaldocument[defos-]{defos}
\externaldocument[cotangent-]{cotangent}
\externaldocument[examples-defos-]{examples-defos}
\externaldocument[algebraic-]{algebraic}
\externaldocument[examples-stacks-]{examples-stacks}
\externaldocument[stacks-sheaves-]{stacks-sheaves}
\externaldocument[criteria-]{criteria}
\externaldocument[artin-]{artin}
\externaldocument[quot-]{quot}
\externaldocument[stacks-properties-]{stacks-properties}
\externaldocument[stacks-morphisms-]{stacks-morphisms}
\externaldocument[stacks-limits-]{stacks-limits}
\externaldocument[stacks-cohomology-]{stacks-cohomology}
\externaldocument[stacks-perfect-]{stacks-perfect}
\externaldocument[stacks-introduction-]{stacks-introduction}
\externaldocument[stacks-more-morphisms-]{stacks-more-morphisms}
\externaldocument[stacks-geometry-]{stacks-geometry}
\externaldocument[moduli-]{moduli}
\externaldocument[moduli-curves-]{moduli-curves}
\externaldocument[examples-]{examples}
\externaldocument[exercises-]{exercises}
\externaldocument[guide-]{guide}
\externaldocument[desirables-]{desirables}
\externaldocument[coding-]{coding}
\externaldocument[obsolete-]{obsolete}
\externaldocument[fdl-]{fdl}
\externaldocument[index-]{index}

% Theorem environments.
%
\theoremstyle{plain}
\newtheorem{theorem}[subsection]{Theorem}
\newtheorem{proposition}[subsection]{Proposition}
\newtheorem{lemma}[subsection]{Lemma}

\theoremstyle{definition}
\newtheorem{definition}[subsection]{Definition}
\newtheorem{example}[subsection]{Example}
\newtheorem{exercise}[subsection]{Exercise}
\newtheorem{situation}[subsection]{Situation}

\theoremstyle{remark}
\newtheorem{remark}[subsection]{Remark}
\newtheorem{remarks}[subsection]{Remarks}

\numberwithin{equation}{subsection}

% Macros
%
\def\lim{\mathop{\mathrm{lim}}\nolimits}
\def\colim{\mathop{\mathrm{colim}}\nolimits}
\def\Spec{\mathop{\mathrm{Spec}}}
\def\Hom{\mathop{\mathrm{Hom}}\nolimits}
\def\Ext{\mathop{\mathrm{Ext}}\nolimits}
\def\SheafHom{\mathop{\mathcal{H}\!\mathit{om}}\nolimits}
\def\SheafExt{\mathop{\mathcal{E}\!\mathit{xt}}\nolimits}
\def\Sch{\mathit{Sch}}
\def\Mor{\mathop{\mathrm{Mor}}\nolimits}
\def\Ob{\mathop{\mathrm{Ob}}\nolimits}
\def\Sh{\mathop{\mathit{Sh}}\nolimits}
\def\NL{\mathop{N\!L}\nolimits}
\def\CH{\mathop{\mathrm{CH}}\nolimits}
\def\proetale{{pro\text{-}\acute{e}tale}}
\def\etale{{\acute{e}tale}}
\def\QCoh{\mathit{QCoh}}
\def\Ker{\mathop{\mathrm{Ker}}}
\def\Im{\mathop{\mathrm{Im}}}
\def\Coker{\mathop{\mathrm{Coker}}}
\def\Coim{\mathop{\mathrm{Coim}}}

% Boxtimes
%
\DeclareMathSymbol{\boxtimes}{\mathbin}{AMSa}{"02}

%
% Macros for moduli stacks/spaces
%
\def\QCohstack{\mathcal{QC}\!\mathit{oh}}
\def\Cohstack{\mathcal{C}\!\mathit{oh}}
\def\Spacesstack{\mathcal{S}\!\mathit{paces}}
\def\Quotfunctor{\mathrm{Quot}}
\def\Hilbfunctor{\mathrm{Hilb}}
\def\Curvesstack{\mathcal{C}\!\mathit{urves}}
\def\Polarizedstack{\mathcal{P}\!\mathit{olarized}}
\def\Complexesstack{\mathcal{C}\!\mathit{omplexes}}
% \Pic is the operator that assigns to X its picard group, usage \Pic(X)
% \Picardstack_{X/B} denotes the Picard stack of X over B
% \Picardfunctor_{X/B} denotes the Picard functor of X over B
\def\Pic{\mathop{\mathrm{Pic}}\nolimits}
\def\Picardstack{\mathcal{P}\!\mathit{ic}}
\def\Picardfunctor{\mathrm{Pic}}
\def\Deformationcategory{\mathcal{D}\!\mathit{ef}}
\usepackage{cleveref}
\crefname{originalreference}{original reference}{original references}
\Crefname{originalreference}{Original reference}{Original references}
\makeatletter
\newlabel{lemma-augmentation-site}{{0D70}{source}{}{sourcepointer-1}{}}
\newlabel{lemma-augmentation-site@cref}{{[originalreference][0][]0D70}{[1][0][]source}{}{}{}}
\makeatother
\crefname{originalreference}{original reference}{original references}
\Crefname{originalreference}{Original reference}{Original references}
\makeatletter
\newlabel{lemma-compare-simplicial-objects-ph-etale}{{0DHJ}{source}{}{sourcepointer-2}{}}
\newlabel{lemma-compare-simplicial-objects-ph-etale@cref}{{[originalreference][0][]0DHJ}{[1][0][]source}{}{}{}}
\makeatother
\crefname{originalreference}{original reference}{original references}
\Crefname{originalreference}{Original reference}{Original references}
\makeatletter
\newlabel{lemma-hypercovering-X-simple-descent-sheaves}{{0DA8}{source}{}{sourcepointer-3}{}}
\newlabel{lemma-hypercovering-X-simple-descent-sheaves@cref}{{[originalreference][0][]0DA8}{[1][0][]source}{}{}{}}
\makeatother
\crefname{originalreference}{original reference}{original references}
\Crefname{originalreference}{Original reference}{Original references}
\makeatletter
\newlabel{lemma-morphism-simplicial-sites}{{0D96}{source}{}{sourcepointer-4}{}}
\newlabel{lemma-morphism-simplicial-sites@cref}{{[originalreference][0][]0D96}{[1][0][]source}{}{}{}}
\makeatother
\crefname{originalsection}{original section}{original sections}
\Crefname{originalsection}{Original section}{Original sections}
\makeatletter
\newlabel{section-hypercovering}{{09A0}{source}{}{sourcepointer-5}{}}
\newlabel{section-hypercovering@cref}{{[originalsection][0][]09A0}{[1][0][]source}{}{}{}}
\makeatother
\crefname{originalsection}{original section}{original sections}
\Crefname{originalsection}{Original section}{Original sections}
\makeatletter
\newlabel{section-simplicial-algebraic-spaces}{{0DE7}{source}{}{sourcepointer-6}{}}
\newlabel{section-simplicial-algebraic-spaces@cref}{{[originalsection][0][]0DE7}{[1][0][]source}{}{}{}}
\makeatother
\crefname{originallemma}{original lemma}{original lemmas}
\Crefname{originallemma}{Original lemma}{Original lemmas}
\makeatletter
\newlabel{spaces-more-cohomology-lemma-comparison-ph-etale}{{0DGT}{source}{}{sourcepointer-7}{}}
\newlabel{spaces-more-cohomology-lemma-comparison-ph-etale@cref}{{[originallemma][0][]0DGT}{[1][0][]source}{}{}{}}
\makeatother
\crefname{originallemma}{original lemma}{original lemmas}
\Crefname{originallemma}{Original lemma}{Original lemmas}
\makeatletter
\newlabel{spaces-more-cohomology-lemma-proper-push-pull-ph-etale}{{0DGX}{source}{}{sourcepointer-8}{}}
\newlabel{spaces-more-cohomology-lemma-proper-push-pull-ph-etale@cref}{{[originallemma][0][]0DGX}{[1][0][]source}{}{}{}}
\makeatother
\crefname{originalsection}{original section}{original sections}
\Crefname{originalsection}{Original section}{Original sections}
\makeatletter
\newlabel{spaces-more-cohomology-section-ph-etale}{{0DGS}{source}{}{sourcepointer-9}{}}
\newlabel{spaces-more-cohomology-section-ph-etale@cref}{{[originalsection][0][]0DGS}{[1][0][]source}{}{}{}}
\makeatother
\crefname{originallemma}{original lemma}{original lemmas}
\Crefname{originallemma}{Original lemma}{Original lemmas}
\makeatletter
\newlabel{spaces-topologies-lemma-surjective-proper-ph}{{0DFI}{source}{}{sourcepointer-10}{}}
\newlabel{spaces-topologies-lemma-surjective-proper-ph@cref}{{[originallemma][0][]0DFI}{[1][0][]source}{}{}{}}
\makeatother
\title{Proper hypercoverings: a\textasciicircum{}\{-1\} on etale sheaves is fully faithful with essential image the cartesian sheaves}
\author{The Stacks Project Authors}
\date{stacks-project:spaces-simplicial.tex:local-snapshot-20261012; GNU FDL 1.2 or later}
\begin{document}
\maketitle
\noindent\textit{Editorial note.} All delivered mathematical text is the author's own, copied from the pinned Stacks Project chapter file sources/stacks-spaces-simplicial/pinned-original-source.tex; no statement or proof step was rewritten, reordered or removed.\par\medskip
\noindent\textit{Editorial note.} The counted proof is the author's complete argument: the fully-faithful half is obtained by what the author calls 'a formal argument (chasing around the diagram)' (source line 8065) through the comparison diagram of lemma-compare-simplicial-objects-ph-etale - all inputs (full faithfulness of h\textasciicircum{}\{-1\} and of a\_\{ph\}\textasciicircum{}\{-1\} with quasi-inverses, the commutativities established in the cited lemmas) are in place before the sentence, and the step is glue between them; the essential-image half is carried by the complete six-equality computation, each equality justified by a named lemma in-span.\par\medskip
\noindent\textit{Editorial note.} The label lemma-descent-sheaves-for-ph-hypercovering names the counted statement, and the proof environment immediately following it in the source is delivered verbatim as its proof.\par\medskip
\noindent\textit{Editorial note.} The counted statement is carried by the author's own lemma environment in the chapter source.\par\medskip
\noindent\textit{Allowed standard prerequisites (editorial).}\begin{itemize}
\item Comparison of the small etale and big ph simplicial sites. Same chapter, \textbackslash{}label\{lemma-compare-simplicial-objects-ph-etale\}. Lines 8046-8048: the diagram h with h\_\{-1\} = a\_X and components h\_n; the fully-faithfulness transfer chases this diagram.
\item The ph-etale comparison lemma. More on Cohomology of Spaces chapter, \textbackslash{}label\{spaces-more-cohomology-lemma-comparison-ph-etale\}. Lines 8049-8052: (h\_\{-1\})\textasciicircum{}\{-1\} is fully faithful with quasi-inverse h\_\{-1,*\}; the same holds for the components h\_n.
\item Description of the simplicial-site functor pair h\textasciicircum{}\{-1\}, h\_*. Same chapter, \textbackslash{}label\{lemma-morphism-simplicial-sites\}. Lines 8053-8054: componentwise full faithfulness lifts to h\textasciicircum{}\{-1\} fully faithful with quasi-inverse h\_*.
\item Simple sheaf descent for hypercoverings on the big ph site. Same chapter, \textbackslash{}label\{lemma-hypercovering-X-simple-descent-sheaves\}. Lines 8061-8064: a\_\{ph\}\textasciicircum{}\{-1\} fully faithful with quasi-inverse a\_\{ph,*\} and essential image the cartesian sheaves on the big total site.
\item Proper pushforward and pullback along the ph-augmentation. More on Cohomology of Spaces chapter, \textbackslash{}label\{spaces-more-cohomology-lemma-proper-push-pull-ph-etale\}. Lines 8097-8102: the third equality of the equalizer computation, using that a\_0 and a\_1 are proper.
\item The augmentation site computes a\_* as an equalizer. Same chapter, \textbackslash{}label\{lemma-augmentation-site\}. Lines 8085-8096: the first and fifth equalities of the computation.
\item Surjective proper morphisms are ph coverings. Topologies on Spaces chapter, \textbackslash{}label\{spaces-topologies-lemma-surjective-proper-ph\}. Lines 8058-8061: U is a hypercovering of X in (Spaces/S)\_ph.
\end{itemize}\medskip
\ifdefined\hypertarget\hypertarget{sourcepointer-1}{}\fi
\noindent\textit{Source reference.} Original reference 0D70: sources/stacks-spaces-simplicial/pinned-original-source.tex, source lines 1141--1141. Prerequisite pointer invoked at proof line(s) 8117; it names a standard fact of the same chapter that the counted argument uses, and every step of the counted proof built on it is retained verbatim in the delivered body.\par\medskip
\ifdefined\hypertarget\hypertarget{sourcepointer-2}{}\fi
\noindent\textit{Source reference.} Original reference 0DHJ: sources/stacks-spaces-simplicial/pinned-original-source.tex, source lines 7973--7973. Prerequisite pointer invoked at proof line(s) 8046; it names a standard fact of the same chapter that the counted argument uses, and every step of the counted proof built on it is retained verbatim in the delivered body.\par\medskip
\ifdefined\hypertarget\hypertarget{sourcepointer-3}{}\fi
\noindent\textit{Source reference.} Original reference 0DA8: sources/stacks-spaces-simplicial/pinned-original-source.tex, source lines 5065--5065. Prerequisite pointer invoked at proof line(s) 8063; it names a standard fact of the same chapter that the counted argument uses, and every step of the counted proof built on it is retained verbatim in the delivered body.\par\medskip
\ifdefined\hypertarget\hypertarget{sourcepointer-4}{}\fi
\noindent\textit{Source reference.} Original reference 0D96: sources/stacks-spaces-simplicial/pinned-original-source.tex, source lines 1273--1273. Prerequisite pointer invoked at proof line(s) 8057; it names a standard fact of the same chapter that the counted argument uses, and every step of the counted proof built on it is retained verbatim in the delivered body.\par\medskip
\ifdefined\hypertarget\hypertarget{sourcepointer-5}{}\fi
\noindent\textit{Source reference.} Original section 09A0: sources/stacks-spaces-simplicial/pinned-original-source.tex, source lines 5012--5012. Prerequisite pointer invoked at proof line(s) 8060; it names a standard fact of the same chapter that the counted argument uses, and every step of the counted proof built on it is retained verbatim in the delivered body.\par\medskip
\ifdefined\hypertarget\hypertarget{sourcepointer-6}{}\fi
\noindent\textit{Source reference.} Original section 0DE7: sources/stacks-spaces-simplicial/pinned-original-source.tex, source lines 6985--6985. Prerequisite pointer invoked at proof line(s) 8039; it names a standard fact of the same chapter that the counted argument uses, and every step of the counted proof built on it is retained verbatim in the delivered body.\par\medskip
\ifdefined\hypertarget\hypertarget{sourcepointer-7}{}\fi
\noindent\textit{Source reference.} Original lemma 0DGT: sources/stacks-chow/tags.txt, source lines 16883--16883. Prerequisite pointer invoked at proof line(s) 8053; it names a standard fact of a sibling chapter that the counted argument uses, and every step of the counted proof built on it is retained verbatim in the delivered body.\par\medskip
\ifdefined\hypertarget\hypertarget{sourcepointer-8}{}\fi
\noindent\textit{Source reference.} Original lemma 0DGX: sources/stacks-chow/tags.txt, source lines 16887--16887. Prerequisite pointer invoked at proof line(s) 8121; it names a standard fact of a sibling chapter that the counted argument uses, and every step of the counted proof built on it is retained verbatim in the delivered body.\par\medskip
\ifdefined\hypertarget\hypertarget{sourcepointer-9}{}\fi
\noindent\textit{Source reference.} Original section 0DGS: sources/stacks-chow/tags.txt, source lines 16882--16882. Prerequisite pointer invoked at proof line(s) 8050; it names a standard fact of a sibling chapter that the counted argument uses, and every step of the counted proof built on it is retained verbatim in the delivered body.\par\medskip
\ifdefined\hypertarget\hypertarget{sourcepointer-10}{}\fi
\noindent\textit{Source reference.} Original lemma 0DFI: sources/stacks-chow/tags.txt, source lines 16838--16838. Prerequisite pointer invoked at proof line(s) 8062; it names a standard fact of a sibling chapter that the counted argument uses, and every step of the counted proof built on it is retained verbatim in the delivered body.\par\medskip
\begin{lemma}
\label{lemma-descent-sheaves-for-ph-hypercovering}
Let $S$ be a scheme. Let $X$ be an algebraic space over $S$.
Let $U$ be a simplicial algebraic space over $S$. Let $a : U \to X$
be an augmentation. If $a : U \to X$ is a proper hypercovering of $X$,
then
$$
a^{-1} : \Sh(X_\etale) \to \Sh(U_\etale)
\quad\text{and}\quad
a^{-1} : \textit{Ab}(X_\etale) \to \textit{Ab}(U_\etale)
$$
are fully faithful with essential image the cartesian sheaves and
quasi-inverse given by $a_*$. Here $a : \Sh(U_\etale) \to \Sh(X_\etale)$
is as in Section \ref{section-simplicial-algebraic-spaces}.
\end{lemma}

\begin{proof}
We will prove the statement for sheaves of sets. It will be an
almost formal consequence of results already established.
Consider the diagram of
Lemma \ref{lemma-compare-simplicial-objects-ph-etale}.
In the proof of this lemma we have seen that
$h_{-1}$ is the morphism $a_X$ of
More on Cohomology of Spaces, Section
\ref{spaces-more-cohomology-section-ph-etale}.
Thus it follows from
More on Cohomology of Spaces, Lemma
\ref{spaces-more-cohomology-lemma-comparison-ph-etale}
that $(h_{-1})^{-1}$ is fully faithful with quasi-inverse $h_{-1, *}$.
The same holds true for the components $h_n$ of $h$.
By the description of the functors $h^{-1}$ and $h_*$ of
Lemma \ref{lemma-morphism-simplicial-sites}
we conclude that $h^{-1}$ is fully faithful with quasi-inverse $h_*$.
Observe that $U$ is a hypercovering of $X$ in $(\textit{Spaces}/S)_{ph}$
as defined in Section \ref{section-hypercovering} since a surjective
proper morphism gives a ph covering by Topologies on Spaces, Lemma
\ref{spaces-topologies-lemma-surjective-proper-ph}.
By Lemma \ref{lemma-hypercovering-X-simple-descent-sheaves}
we see that $a_{ph}^{-1}$ is fully faithful with quasi-inverse
$a_{ph, *}$ and with essential image the cartesian sheaves
on $(\textit{Spaces}/U)_{ph, total}$.
A formal argument (chasing around the diagram) now shows that
$a^{-1}$ is fully faithful.

\medskip\noindent
Finally, suppose that $\mathcal{G}$ is a cartesian sheaf on $U_\etale$.
Then $h^{-1}\mathcal{G}$ is a cartesian sheaf on
$(\textit{Spaces}/U)_{ph, total}$.
Hence $h^{-1}\mathcal{G} = a_{ph}^{-1}\mathcal{H}$ for some sheaf
$\mathcal{H}$ on $(\textit{Spaces}/X)_{ph}$.
We compute using somewhat pedantic notation
\begin{align*}
(h_{-1})^{-1}(a_*\mathcal{G})
& =
(h_{-1})^{-1}
\text{Eq}(
\xymatrix{
a_{0, small, *}\mathcal{G}_0
\ar@<1ex>[r] \ar@<-1ex>[r] &
a_{1, small, *}\mathcal{G}_1
}
) \\
& =
\text{Eq}(
\xymatrix{
(h_{-1})^{-1}a_{0, small, *}\mathcal{G}_0
\ar@<1ex>[r] \ar@<-1ex>[r] &
(h_{-1})^{-1}a_{1, small, *}\mathcal{G}_1
}
) \\
& =
\text{Eq}(
\xymatrix{
a_{0, big, ph, *}h_0^{-1}\mathcal{G}_0
\ar@<1ex>[r] \ar@<-1ex>[r] &
a_{1, big, ph, *}h_1^{-1}\mathcal{G}_1
}
) \\
& =
\text{Eq}(
\xymatrix{
a_{0, big, ph, *}(a_{0, big, ph})^{-1}\mathcal{H}
\ar@<1ex>[r] \ar@<-1ex>[r] &
a_{1, big, ph, *}(a_{1, big, ph})^{-1}\mathcal{H}
}
) \\
& =
a_{ph, *}a_{ph}^{-1}\mathcal{H} \\
& =
\mathcal{H}
\end{align*}
Here the first equality follows from Lemma \ref{lemma-augmentation-site},
the second equality follows as $(h_{-1})^{-1}$ is an exact functor,
the third equality follows from
More on Cohomology of Spaces, Lemma
\ref{spaces-more-cohomology-lemma-proper-push-pull-ph-etale}
(here we use that $a_0 : U_0 \to X$ and $a_1: U_1 \to X$ are proper),
the fourth follows from $a_{ph}^{-1}\mathcal{H} = h^{-1}\mathcal{G}$,
the fifth from Lemma \ref{lemma-augmentation-site}, and the
sixth we've seen above. Since $a_{ph}^{-1}\mathcal{H} = h^{-1}\mathcal{G}$
we deduce that $h^{-1}\mathcal{G} \cong h^{-1}a^{-1}a_*\mathcal{G}$
which ends the proof by fully faithfulness of $h^{-1}$.
\end{proof}

\end{document}
